\documentclass{article}
\usepackage[T1]{fontenc}
\usepackage{lmodern}
\usepackage{graphicx}
\usepackage{amsmath,amssymb}
\usepackage[a4paper,margin=2cm]{geometry}
\usepackage[absolute,overlay]{textpos}
\setlength{\TPHorizModule}{1mm}
\setlength{\TPVertModule}{1mm}
\usepackage[indonesian]{babel}
\usepackage{hyperref}
\setlength{\emergencystretch}{2em}
% unit: Halaman 48

\begin{document}

% === Halaman 48 (python/native) ===

Sejarah Kriptografi

• Kriptografi sudah berusia sangat tua, sudah ada sejak peradaban manusia di bumi

• Secara historis, kriptografi diasosiasikan dengan kegiatan mata-mata, 

pemerintahan, dan militer, dan telah digunakan di dalam perang selama ribuan 

tahun.

• Tiga pihak yang memiliki kontribusi penting di dalam perkembangan kriptografi 

zaman dahulu: kalangan militer, diplomat, dan diarist.

• Sejak lebih dari 50 tahun yang lalu, kriptografi mendapatkan landasan matematika, 

dan telah bergeser dari aplikasi militer ke aplikasi komersil. 

• Secara garis besar, kriptografi dibagi menjadi dua era: kriptografi klasik dan 

kriptografi modern.

48

\end{document}
